% -*- coding: utf-8 -*-
% !TEX program = xelatex
\documentclass[plain,noanswer,medmath]{../../template/jnuexam}
\usepackage{../../template/kysx}

\begin{document}


\examtitle{2002}{4}


\loadproblems{1}{2002P1}%载入试卷一
\loadproblems{2}{2002P2}%载入试卷二
\loadproblems{3}{2002P3}%载入试卷三


\makepart{填空题}{本题共5小题，每小题3分，满分15分}


\useproblem{3}{1}{1}%【同试卷三第一(1)题】


\begin{problem}
已知$f(x)$的一个原函数为$\ln^2 x$，则$\int xf'(x)\dx =$ \fillin{}．
\end{problem}


\begin{problem}
设矩阵$A = \begin{pmatrix}
  1 & -1 \\
  2 & 3
\end{pmatrix}$，$B = A^2 - 3A + 2E$，则$B^{-1} =$ \fillin{}．
\end{problem}


\begin{problem}
设向量组$\alpha_1 = (a,0,c)$,$\alpha_2 = (b,c,0)$,$\alpha_3 = (0,a,b)$线性无关，
则$a,b,c$必满足关系式 \fillin{}．
\end{problem}

%类似试卷三第一(4)题
\begin{problem}
设随机变量$X$和$Y$的联合概率分布为
$$\begin{array}{|c|c|c|c|}
\hline
  \diagboxtwo{X}{Y} & -1 & 0 & 1 \\
\hline
  0 & 0.07 & 0.18 & 0.15 \\
\hline
  1 & 0.08 & 0.32 & 0.20 \\
\hline
\end{array}$$
则$X$和$Y$的相关系数$\rho=$．
\end{problem}


\makepart{选择题}{本题共5小题，每小题3分，共15分}


\useproblem{3}{2}{1}%【同试卷三第二(1)题】


\useproblem{2}{2}{2}%【同试卷二第二(2)题】


\begin{problem}
设$A$,$B$为$n$阶矩阵，$A^*$,$B^*$分别为$A$,$B$对应的伴随矩阵，
分块矩阵$C = \begin{pmatrix}
  A&O \\
  O&B
\end{pmatrix}$，则$C$的伴随矩阵$C^* = $\pickout{}
\begin{abcd}
\item $\begin{pmatrix}
  |A|A^*  & O \\
  O       & |B|B^* 
\end{pmatrix}$．
\item $\begin{pmatrix}
  |B|B^* & O \\
  O      & |A|A^*
\end{pmatrix}$．
\item $\begin{pmatrix}
  |A|B^* & O \\
  O      & |B|A^*
\end{pmatrix}$．
\item $\begin{pmatrix}
  |B|A^* & O \\
  O      & |A|B^*
\end{pmatrix}$．
\end{abcd}
\end{problem}


\useproblem{1}{2}{5}%【同试卷一第二(5)题】


\begin{problem}
设随机变量$X_1$, $X_2$, $\cdots$, $X_n$相互独立，$S_n = X_1 + X_2 +  \cdots  + X_n$，
则根据列维—林德伯格(Levy-Lindberg)中心极限定理，当$n$充分大时，$S_n$近似服从正态分布，
只要$X_1$, $X_2$, $\cdots$, $X_n$\pickout{}
\begin{abcd}
\item 有相同的数学期望．
\item 有相同的方差．
\item 服从同一指数分布．
\item 服从同一离散型分布．
\end{abcd}
\end{problem}


\makepart{}{本题满分5分}%【同试卷三第三题】

\useproblem{3}{3}{0}


\makepart{}{本题满分7分}%【同试卷三第四题】

\useproblem{3}{4}{0}


\makepart{}{本题满分6分}%【同试卷三第五题】

\useproblem{3}{5}{0}


\makepart{}{本题满分7分}

\begin{problem}
设闭区域$D:{x^2} + {y^2} \leqslant y,x \geqslant 0$；$f(x,y)$为$D$上的连续函数，且
$$f(x,y) = \sqrt {1 - {x^2} - {y^2}}  - \frac{\pi }{8}\iint {f(u,v)\du\dv},$$
求$f(x,y)$．
\end{problem}


\makepart{}{本题满分7分}

\begin{problem}
设某商品需求量$Q$是价格$p$的单调减少函数：$Q = Q(p)$，
其需求弹性$\eta  = \frac{{2{p^2}}}{{192 - {p^2}}} > 0$．
\begin{enumerate}
\item 设$R$为总收益函数，证明$\frac{{\d R}}{{\d p}} = Q(1 - \eta )$．
\item 求$p = 6$时，总收益对价格的弹性，并说明其经济意义．
\end{enumerate}
\end{problem}


\makepart{}{本题满分6分}%【同试卷三第八题】

\useproblem{3}{8}{0}


\makepart{}{本题满分8分}

\begin{problem}
设四元齐次线性方程组①为$\begin{cases}
  2x_1 + 3x_2 - x_3 = 0 \\
  x_1 + 2x_2 + x_3 - x_4 = 0
\end{cases}$，且已知另一四元齐次线性方程组②的一个基础解系为
$\alpha_1 = (2, - 1,a + 2,1)^T$, $\alpha_2 = (- 1,2,4,a + 8)^T$．
\begin{enumerate}
\item 求方程组①的一个基础解系；
\item 当$a$为何值时，方程组①与②有非零公共解？在有非零公共解时，求出全部非零公共解．
\end{enumerate}
\end{problem}


\makepart{}{本题满分8分}

\begin{problem}
设实对称矩阵$A = \begin{pmatrix}
  a &  1 & 1 \\
  1 &  a & -1 \\
  1 & -1 & a
\end{pmatrix}$，求可逆矩阵$P$，使$P^{-1}AP$为对角矩阵，并计算行列式$|A - E|$的值．
\end{problem}


\makepart{}{本题满分8分}

\begin{problem}
设$A$、$B$是任意二事件，其中$A$的概率不等于$0$和$1$，
证明：$P(B|A) = P(B|\bar A)$是事件$A$与$B$独立的充分必要条件．
\end{problem}


\makepart{}{本题满分8分}%【同试卷三第十二题】

\useproblem{3}{12}{0}


\end{document}
